% Generated draft: requires mathematical and source review.
\documentclass[11pt]{article}
\usepackage[margin=25mm]{geometry}
\usepackage{fontspec}
\setmainfont{Latin Modern Roman}
\usepackage{amsmath,amssymb,mathtools,unicode-math}
\setmathfont{Latin Modern Math}
\usepackage{enumitem,needspace,xurl,hyperref}
\setlength{\parindent}{0pt}
\setlength{\parskip}{5pt}
\setlength{\emergencystretch}{3em}
\newcommand{\N}{\mathbb{N}}
\newcommand{\Z}{\mathbb{Z}}
\newcommand{\Q}{\mathbb{Q}}
\newcommand{\R}{\mathbb{R}}
\newcommand{\C}{\mathbb{C}}
\newcommand{\F}{\mathbb{F}}
\newcommand{\A}{\mathbb{A}}
\newcommand{\PP}{\mathbb{P}}
\newcommand{\E}{\mathbb{E}}
\newcommand{\eps}{\varepsilon}
\newcommand{\dd}{\,\mathrm{d}}
\begin{document}
\section*{UnsolvedMath ID: 701590}
\textbf{Source problem code:} catalogue entry
\subsection{Definitions and mathematical statement}
A group is word-hyperbolic if it admits a finite generating set whose Cayley graph, with its graph metric, has uniformly thin geodesic triangles. For a word-hyperbolic group \(K\), write \(\partial K\) for its Gromov boundary, defined by equivalence classes of geodesic rays, where two rays are equivalent if their images have finite Hausdorff distance. Equip \(K\cup\partial K\) with its Gromov compactification topology.

Let \(G\) be a word-hyperbolic group and let \(H\leq G\) be a subgroup that is word-hyperbolic with respect to its own word metric. Suppose the inclusion \(i:H\hookrightarrow G\) extends to a continuous map
\[
\overline{i}:H\cup\partial H\longrightarrow G\cup\partial G
\]
whose boundary restriction \(j=\overline{i}|_{\partial H}\) takes values in \(\partial G\) and satisfies
\[
j(h\cdot\xi)=h\cdot j(\xi)\qquad(h\in H,\ \xi\in\partial H).
\]
The boundary extension, when it exists, is unique and is called the Cannon--Thurston map. For \(p\in\partial G\), its full fiber is
\[
j^{-1}(p)=\{\xi\in\partial H:j(\xi)=p\}.
\]
For every pair \((G,H)\) satisfying these hypotheses, ask:
\begin{enumerate}
\item[(a)] Must every fiber be finite, that is,
\[
\forall p\in\partial G,\qquad |j^{-1}(p)|<\infty?
\]
\item[(b)] Must there exist a positive integer \(N=N(G,H)\) such that
\[
\forall p\in\partial G,\qquad |j^{-1}(p)|\leq N?
\]
\end{enumerate}
No surjectivity assumption is imposed in the general questions; points outside the image have empty fibers. The source imposes no additional non-elementarity hypothesis.

As background, the supplied source reports affirmative answers to both questions in all cases in which the Cannon--Thurston map was known to exist and its multiplicity had been analyzed. It includes the original Cannon--Thurston setting: \(G=\pi_1(M)\), where \(M\) is a closed hyperbolic \(3\)-manifold fibering over the circle with surface fiber \(S\), and \(H\) is the subgroup corresponding to \(\pi_1(S)\). In that setting the source describes
\[
j:\mathbb{S}^1\longrightarrow\mathbb{S}^2
\]
as a continuous, surjective, uniformly finite-to-one ``Peano curve.'' These are supplied contextual assertions, not an independently verified assessment of the literature.
\subsection{Short English statement}
For every word-hyperbolic subgroup \(H\) of a word-hyperbolic group \(G\) whose inclusion admits a Cannon--Thurston boundary map, must each fiber be finite? Must the fiber sizes also have a finite bound depending only on \((G,H)\)?
\subsection{Sources}
\begin{itemize}
\item Local UnsolvedMath registry, permanent ID \(701590\), source problem code \texttt{catalogue entry}, titled ``Kourovka Notebook Problem 19.42.'' The supplied record labels the problem open; this label is not independently verified.
\item I. Kapovich, proposer, 2018. The supplied attribution identifies the original problem as Problem 19.42 in \emph{The Kourovka Notebook (Unsolved Problems in Group Theory)}, 19th issue (2018), and also cites the 21st issue edition (2026), \texttt{arXiv:1401.0300v48}.
\item The supplied collection description identifies E. I. Khukhro and V. D. Mazurov as editors of the Kourovka Notebook collection covering Issues 1--21 (1965--2026). The original Cannon--Thurston example and the statement about analyzed cases are supplied background; no separate bibliographic citation for that example is provided.
\end{itemize}
\end{document}
